\section{Introduction}
\IEEEPARstart{D}{eep} slow-wave sleep (N3) is one of the most important states the brain enters each night. Its cortical slow oscillations are strongest early in the night and weaken as sleep pressure falls \cite{borbely2016twoprocess}, and the amplitude and slope of single slow waves decline with them \cite{riedner2007slowwaves,tononi2014plasticity}. In animals, sleep increases the removal of waste products such as $\beta$-amyloid from the brain \cite{xie2013clearance}. In community cohorts, less slow-wave sleep goes with more signs of brain ageing \cite{baril2021brainaging} and, over several years, with a higher dementia risk \cite{himali2023dementia}. Slow-wave measures are therefore promising markers for brain-health research, especially over many nights at home.

Three problems stand out. First, clinical scoring turns a graded process into one label: an epoch is N3 when at least 20\% of it has 0.5--2~Hz waves of at least 75~\textmu V peak-to-peak \cite{aasm2020manual}, whereas the older Rechtschaffen and Kales rules split this range into stage 3 (20--50\%) and stage 4 (more than 50\%) \cite{rechtschaffen1968manual}. An epoch that just passes the threshold and one full of large waves therefore get the same score. Second, the N2/N3 boundary is the least reliable decision in manual scoring: N3 agreement was only 67.4\% across more than 2,500 scorers in the AASM inter-scorer programme \cite{rosenberg2013interscorer}, and the fixed 75~\textmu V criterion is criticised because EEG amplitude varies with age and sex \cite{davidson2025n3,rosinvil2021age}. Third, N3 duration varies far more from night to night than slow-wave activity does \cite{qin2026swa,mullins2024stability}, so a count of N3 epochs is a rough summary of the physiology that matters most.

Single-channel EEG makes long-term recording practical, and deep learning stagers now reach expert-level agreement on such data \cite{supratak2017deepsleepnet,supratak2020tinysleepnet,eldele2021attnsleep,phan2023lseqsleepnet,perslev2021usleep,vallat2021yasa,phan2022review}. Their output, however, is still a binary label or a stage probability. Continuous depth measures such as the odds ratio product \cite{younes2015orp} and the sleep depth index \cite{zhou2025sdi} describe depth across all stages, not the strength of slow waves inside N3. What is missing is a single-channel method that finds N3, describes slow-wave strength in each epoch and night in clear physiological terms, and warns when its output should not be trusted.

This work presents such a method and tests it as strictly as public data allow. The detector learns from expert labels; the description uses fixed rules, because no dataset contains expert ratings of slow-wave strength and a model trained on its own input features would only confirm itself. We fixed seven research questions with two sub-questions and their success criteria before computing the full results, and report every verdict as computed. Across seven cohorts and 457 participants, the detector ranked N3 epochs almost as well in independent centres and devices as in cross-validation and matched or outperformed four published architectures retrained under the same protocol; the factor scores turned the binary label into a graded description that follows the night-time decline, age and the graded expert vote; and the reliability score flagged likely errors and, within ISRUC-I, the epochs on which the two experts disagreed.

The main contributions are:
\begin{enumerate}
\item A small single-channel N3 detector (96,226 parameters) with a dilated temporal-context block, evaluated with participant-grouped cross-validation on 251 participants and on five external cohorts, and compared with four published architectures retrained under the same protocol and with a pretrained open-source stager. % rq1_detection_metrics
\item Four rule-based factor scores (delta intensity, slow-wave morphology, spectral dominance and N3 continuity), an exactly additive overall rating, and a Low/Moderate/High slow-wave grade based on the published 20\% and 50\% coverage criteria.
\item An evaluation of the description against the homeostatic decline, the graded vote of two to five experts per epoch, expert Rechtschaffen and Kales stage 3 versus stage 4, age, and, descriptively, night-to-night stability.
\item A reliability score that keeps signal quality, confidence and out-of-distribution distance as separate inputs; removing the 20\% least reliable epochs raised F1 from 0.849 to 0.949 on the remaining epochs. % rq6_risk_coverage
\end{enumerate}
